\documentclass[aspectratio=169,10pt]{beamer}

\usetheme{Boadilla}
\usecolortheme{default}
\usepackage[T1]{fontenc}
\usepackage[utf8]{inputenc}
\usepackage{lmodern}

\definecolor{TBLNavy}{RGB}{20,35,58}
\definecolor{TBLBlue}{RGB}{40,111,160}
\definecolor{TBLTeal}{RGB}{24,145,137}
\definecolor{TBLLight}{RGB}{239,245,249}
\setbeamercolor{structure}{fg=TBLBlue}
\setbeamercolor{title}{fg=white,bg=TBLNavy}
\setbeamercolor{frametitle}{fg=white,bg=TBLNavy}
\setbeamercolor{block title}{fg=white,bg=TBLBlue}
\setbeamercolor{block body}{fg=TBLNavy,bg=TBLLight}
\setbeamercolor{alerted text}{fg=TBLTeal}
\setbeamertemplate{navigation symbols}{}
\setbeamertemplate{footline}[frame number]
\setbeamersize{text margin left=10mm,text margin right=10mm}
\setlength{\parskip}{0.5em}
\setlength{\emergencystretch}{3em}

\title[Linux Resource Monitor]{Speaking Notes: Linux System Resource Monitoring}
\subtitle{Five technical sources behind the presentation}
\author{Yiğit Cemal Öztürk \and Murat Karatekin}
\institute{Kocaeli University \\ TBL221 Linux Operating System Usage and Administration}
\date{}

\begin{document}

\begin{frame}[plain]
  \titlepage
\end{frame}

\begin{frame}{1. Linux Kernel Documentation: ``The /proc Filesystem''}
\begin{block}{Speaking note}
The CPU values exposed through \texttt{/proc/stat} are cumulative counters, not a ready-made utilization percentage. Our recorder compares successive samples and derives utilization from their differences, so the sampling interval matters.
\end{block}
\begin{itemize}
  \item Connect this source to the CPU measurement row in the main presentation.
  \item Keep the distinction clear: the kernel exposes counters; the recorder calculates the percentage.
\end{itemize}
\end{frame}

\begin{frame}{2. Linux man-pages: \texttt{proc(5)}}
\begin{block}{Speaking note}
This manual documents process and memory information exposed through \texttt{/proc}. For memory, emphasize \texttt{MemAvailable} rather than interpreting the \texttt{free} field alone; the project's reported used-memory value is total minus available.
\end{block}
\begin{itemize}
  \item Connect this source to process counting and memory measurement.
  \item The project uses \texttt{MemAvailable} as part of its used-memory calculation.
\end{itemize}
\end{frame}

\begin{frame}{3. SQLite: ``Write-Ahead Logging''}
\begin{block}{Speaking note}
WAL lets readers continue using a stable database snapshot while the recorder appends new records. This suits our one-writer, read-only-dashboard design. Do not imply unlimited concurrency or describe WAL as a backup mechanism.
\end{block}
\begin{itemize}
  \item Connect this source to the separation between the recorder and dashboard.
  \item The relevant use case is one writer with a concurrent reader.
\end{itemize}
\end{frame}

\begin{frame}{4. Python Software Foundation: Python 3 Library Reference}
\begin{block}{Speaking note}
The recorder's dependency-free claim means it relies on Python's standard library for its runtime needs. This reference documents those APIs, but it is not evidence of the project's test results. PySide6 remains a separate GUI dependency.
\end{block}
\begin{itemize}
  \item Connect this source to the local-run and implementation slide.
  \item Distinguish the recorder's standard-library-only runtime from the GUI dependency.
\end{itemize}
\end{frame}

\begin{frame}{5. Freedesktop.org: \texttt{systemd.exec}}
\begin{block}{Speaking note}
Options such as \texttt{NoNewPrivileges}, \texttt{ProtectSystem}, and restricted writable paths reduce the service's privileges and filesystem access. Present them as defense in depth, not as a guarantee of complete isolation.
\end{block}
\begin{itemize}
  \item Connect this source to the service and packaging slide.
  \item The aim is to limit the service's access, not to claim that all risks are eliminated.
\end{itemize}
\end{frame}

\end{document}
